\section{侦察终结、侦查羁押期限、补充侦查}
\begin{enumerate}
  \item 侦查终结。
        \begin{enumerate}
          \item 事实清楚；证据确实、充分；法律手续完备。
          \item 辩护人提出要求的应当听他的意见。被动听。
          \item 后续处理。移送同级检察院审查起诉；撤销案件。
        \end{enumerate}
  \item 侦查羁押期限。最长可以逮捕多长时间，\zd{2 + 1 + 2 + 2 + X}。
        \begin{enumerate}
          \item 一般 2 个月。
          \item 案情复杂，上一级检察院，+1 个月。
          \item 交集流广重大复杂案件，省一级检察院，+2 个月。
          \item 10 年以上，省一级检察院，+2 个月。凡是延长 2 个月的都得省一级检察院批准。
          \item 特殊原因，最高检报全国人大常委，+X 月。
          \item 不包括逮捕之前的先行拘留的时间。
        \end{enumerate}
  \item \zd{补充侦查}。4 种常见的理由。\zd{事实不清，证据不足，漏罪漏人}。
        \begin{enumerate}
          \item 检察院决定，不同阶段；
                \begin{enumerate}
                  \item 批捕。说明理由，通知公安补侦。
                  \item 审查起诉。
                        \begin{enumerate}
                          \item 哪里来哪里去，\zd{捕诉部门}自己来；自侦 \hz 监委 退回 > 自己来；
                          \item 每次 1 月，\zd{累计 2 次}，不管管辖转变。
                          \item 退 1/2 次补侦证据不足 \yjt 可以/应当不起诉。
                        \end{enumerate}
                  \item 审判。
                        \begin{enumerate}
                          \item 检察院启动。例外：被告人提供了新的立功线索，\zd{可以建议}补侦 \bdy 可以补侦。
                          \item 补侦主体：原则检察院（已经指控的事实）\hz 公安（漏罪）。
                          \item 期限同上；
                        \end{enumerate}
                \end{enumerate}
        \end{enumerate}
\end{enumerate}
